\section{Eliminating repeated kernel multiplicity in an alternating carrier}
\label{sec:heisenberg-repeated-kernel-exclusion}

\paragraph{Origin and review status.}
D-TD33, reserved question Q2eq, is the root's combined actual-source
argument from the separately checked D-TD27--32 and HD-FT27.
The component derivations and their independent audits retain their
separate attribution. Sixfolds checked the proposed interfaces and
pointed out the useful noncircular ordering of the degree estimate.
The complete author proof and dependency-interface reading passes;
the independent whole-proof review also passes.

\paragraph{Actual source hypotheses.}
Let $J\subset\operatorname{GL}(V)$ be an actual faithful irreducible
finite complex coefficient image of a connected cover $C\to\mathbf P^1$.
Let $T_i$ be its original inertia operators and put $e=\dim V$.
Suppose there is a stable whole tensor factorization $V=W\otimes R$,
with the entire arbitrary partner $R$ retained, such that
\[
 \sum_i\operatorname{rank}(1-T_i)=2e+6,\qquad
 P=\operatorname{im}(J\to\operatorname{PGL}(W)),\qquad
 \mathcal Q\triangleleft P,\quad
 \mathcal Q\simeq\mathbf F_p^{D},\quad D=2m\ge2.
 \tag{TDR.1}
\]
Assume the projective action of $\mathcal Q$ on $W$ is irreducible
with nondegenerate scalar commutator form. As in D-TD05, its actual
linear conjugation image is
$H=P/\mathcal Q\le\operatorname{Sp}(\mathcal Q)$ and
$\dim W=p^{D/2}$. Require that $H$ acts irreducibly on the original
$\mathbf F_p$-space $\mathcal Q$.

\begin{theorem}[D-TD33]
In this setup, suppose there is an actual surjection
\[
 \psi:H\longrightarrow A_n,\qquad N=\ker\psi,\qquad
 n\ge\max\{10^6,B+2\},\qquad B=28900^{14}.
 \tag{TDR.2}
\]
The original-field restriction $\mathcal Q|_N$ cannot be homogeneous
of multiplicity greater than one. Consequently, in the homogeneous
case it is irreducible over $\mathbf F_p$.
\end{theorem}

\begin{proof}
Suppose $\mathcal Q|_N=S^r$ with $S$ simple over $\mathbf F_p$
and $r>1$. Write $d=\dim_{\mathbf F_p}S$, so $D=dr$.
D-TD27 constructs an actual faithful projective action
$A_n\hookrightarrow\operatorname{PGL}_r(F)$ on the multiplicity
space, over the intrinsic finite field $F=\operatorname{End}_N(S)$.
This conclusion precedes the bounded-kernel and central-product
steps of D-TD28--29. D-TD32 therefore already gives
\[
                  r>n/5,\qquad D=dr>n/5,\qquad
                  e\ge\dim W=p^{D/2}>2^{n/10}.
 \tag{TDR.3}
\]
No actual quasisimple subgroup has yet been used in this deduction.

We record one universal numerical comparison for later use.
For $n\ge10^6$, let $q=\lfloor\log_2 n\rfloor\ge19$.
At $q=19$ the inequality $2^{q-12}>3(q+1)$ reads $128>60$.
Its left-minus-right difference increases at the next step by
$2^{q-12}-3>0$, proving it for every larger integer $q$.
Since $2^q\le n<2^{q+1}$, we obtain
\[
       \frac{n}{4096}\ge2^{q-12}>3(q+1)>3\log_2 n,
       \qquad 2^{n/4096}>n^3.
 \tag{TDR.4}
\]
Thus (TDR.3) implies $e>n^3>504$ before the use of D-TD28.

All hypotheses of D-TD28--29 now hold on this same actual source.
They give $|N|<B$ and construct an actual normal quasisimple
subgroup $E\triangleleft H$ with
\[
 H=NE,\qquad [N,E]=1,\qquad N\cap E=Z(E),\qquad
 E/Z(E)\simeq A_n,
 \qquad
 \mathcal Q=F_0^a\otimes_{F_0}F_0^s,
 \tag{TDR.5}
\]
where $F_0$ is a finite extension of $\mathbf F_p$,
$f_0=[F_0:\mathbf F_p]$, and $D=af_0s$.
The second tensor factor is an actual faithful absolutely
irreducible $E$-module over $F_0$, with projective kernel exactly
$Z(E)$. The restriction of $\psi$ identifies its projective
quotient with this same $A_n$. In particular D-TD32 applied again,
now to the absolute factor, gives $s>n/5$.
The original multiplicity $r$ and the absolute degree $s$ are
not silently identified; D-TD29 gives $r=[F_0:F]s$.

Let $h\in H$ have nonidentity image $y=\psi(h)$, moving $k>0$
of the $n$ letters. Write $h=uz$ with $u\in N$ and $z\in E$.
In (TDR.5) its matrix is $A_u\otimes B_z$, and $z$ maps to $y$.
The spectral quantity
$\nu(B_z)=s-\max_\lambda\dim\ker(B_z-\lambda I)$ is computed
after scalar extension to an algebraic closure; ranks are unchanged
by that extension. D-TD31 for this actual perfect central cover
and D-TD27's all-characteristic tensor eigenspace inequality yield
\[
 \begin{split}
 \nu(B_z)&\ge\frac{s k}{24n}>\frac{k}{120},\\
 \operatorname{rank}_{\mathbf F_p}(1-h)
    &=f_0\operatorname{rank}_{F_0}(1-A_u\otimes B_z)
      \ge af_0\nu(B_z)>\frac{k}{120}.
 \end{split}
 \tag{TDR.6}
\]
This uses geometric eigenspaces and remains valid for nonsemisimple
operators in characteristic $p$. It holds for every element of $H$
with nonidentity alternating image, including all such powers of
every original inertia operator.

The exact Heisenberg trace formula D-TD05 consequently gives, for
a unitary representative $A$ on $W$ inducing $h$,
\[
 \frac{|\operatorname{tr}A|}{\dim W}
 \le p^{-\operatorname{rank}(1-h)/2}
 \le 2^{-\operatorname{rank}(1-h)/2}
 <2^{-k/240},\qquad
 \frac{|\operatorname{tr}A|}{\dim W}\le\frac1{\sqrt2}<\frac78.
 \tag{TDR.7}
\]
The formula also allows zero trace. Its last bound follows because
the positive integer rank is at least one.

Let $f(y)=n-k$ denote the number of fixed letters. If
$0<k\le n/16$, then
\[
 \frac{|\operatorname{tr}A|}{\dim W}
 \le\frac1{\sqrt2}\le\left(\frac{15}{16}\right)^3
 \le\left(\frac{f(y)}n\right)^3.
 \tag{TDR.8}
\]
The middle inequality is exact: its square is
$16^6\le2\cdot15^6$, whose positive margin is $6004034$.
If $k>n/16$, (TDR.7) and (TDR.4) give instead
\[
 \frac{|\operatorname{tr}A|}{\dim W}
       <2^{-n/3840}<2^{-n/4096}<\frac1{n^3}.
 \tag{TDR.9}
\]
For the identity image the trivial trace bound is one. Together
these are exactly the cubic fixed-letter estimate for every
element of the actual projective factor.

We verify its transfer to the full original group and every power.
The finite projective group $P$ has a finite inverse image in
$\operatorname{SL}(W)$, whose scalar kernel has order $\dim W$.
Unitarize that finite group once. For any $g\in J$ choose a
representative $A$ there and write its actual operator as
$T_g=A\otimes B$, with the reciprocal scalar absorbed in $B$.
If $L=\operatorname{ord}_J(g)$, then $A^L=\zeta I$ for a root
of unity $\zeta$, and $B^L=\zeta^{-1}I$ because $T_g^L=I$.
Thus $B$ has finite order, and for every integer $j$ its normalized
absolute trace is at most one. The same fixed choice gives
$T_g^j=A^j\otimes B^j$ for all $j$. Hence the entire partner
preserves every bound above, without replacing $L$ by a quotient
order or dropping its kernel powers.

The composite $x:J\to P\to H\xrightarrow{\psi}A_n$ is an
actual surjection. We have proved for every $g\in J$ that
\[
 \frac{|\operatorname{tr}T_g|}{e}
       \le\frac{\max\{f(x(g)),1\}^3}{n^3},\qquad
 x(g)\ne1\ \Longrightarrow
       \frac{|\operatorname{tr}T_g|}{e}\le\frac1{\sqrt2}<\frac78,
 \qquad e>n^3.
 \tag{TDR.10}
\]
All hypotheses of the accepted actual-source cubic obstruction
HD-FT27 are satisfied. Its full-order Fourier and actual triple-genus
argument contradicts (TDR.1). This proves the theorem.
\end{proof}

\paragraph{Application to the maximal alternating layer.}
In the actual faithful-partner $A_n$ branch of D-TD19 at the same
maximal natural degree $n\ge N_*$, D-TD21 supplies the original-field
homogeneity required here. D-TD29 verifies from the explicit
definition of $N_*$ that $N_*>B+2$, and $N_*>10^6$ is immediate
from that same definition. Thus the kernel restriction in this
branch is now irreducible over $\mathbf F_p$.
This conclusion concerns the actual Heisenberg carrier kernel,
not the characteristic-zero kernel of the entire standard partner.

\paragraph{Retained boundary.}
The theorem imposes no maximality when its own hypotheses hold.
It does not treat multiplicity one, an $S_n$ quotient or sign-exchanged
components, and does not infer absolute irreducibility from
irreducibility over the prime field. No general bound on $m$ or $n$,
target source, exact complex field, polarization transport or detecting
cycle is supplied by this exclusion. All numerical comparisons above
are universal written inequalities; no new finite computation is
needed or claimed.
